\documentclass[11pt]{amsart}
\usepackage[T1]{fontenc}
\usepackage[utf8]{inputenc}
\usepackage{amssymb}
\usepackage{array}
\usepackage[margin=1.05in]{geometry}
\usepackage[hidelinks]{hyperref}
\newcommand{\stamptime}{\the\numexpr\time/60\relax:\ifnum\numexpr\time-60*(\time/60)\relax<10 0\fi\the\numexpr\time-60*(\time/60)\relax}
\AddToHook{shipout/foreground}{\put(0.5\paperwidth,-\paperheight+0.55cm){\makebox[0pt]{\tiny\ttfamily unit/doc/\jobname.tex \ --- \ compiled \today, \stamptime}}}
\newtheorem{theorem}{Theorem}
\newtheorem{proposition}[theorem]{Proposition}
\newtheorem{lemma}[theorem]{Lemma}
\newtheorem{corollary}[theorem]{Corollary}
\theoremstyle{definition}
\newtheorem{definition}[theorem]{Definition}
\theoremstyle{remark}
\newtheorem{remark}[theorem]{Remark}
\newcommand{\Om}{\Omega}
\usepackage{url}
% declaration names: typewriter, breakable after . - [ ] (no stretched lines)
\DeclareUrlCommand\urlcode{\urlstyle{tt}}
\DeclareRobustCommand{\code}[1]{\ifmmode\text{\texttt{#1}}\else\urlcode{#1}\fi}
\renewcommand{\UrlBreaks}{\do\.\do\-\do\[\do\]}
\newcommand{\cT}{\mathsf{isTrunc}}
\newcommand{\cC}{\mathsf{isConn}}
\newcommand{\refl}{\mathsf{refl}}
\newcommand{\ap}{\mathsf{ap}}
\newcommand{\Fib}{\mathrm{Fib}}
\newcommand{\fpr}{\mathsf{fpr}}
\newcommand{\SCt}{\mathsf{SC3}}
\newcommand{\SCf}{\mathsf{SC4}}
\newcommand{\Tri}{\mathsf{Triples}}
\newcommand{\Mor}{\mathsf{Mor}}
\newcommand{\Hom}{\mathrm{Hom}}
\newcommand{\Map}{\mathrm{Map}_*}
\newcommand{\real}{\mathsf{realise}}
\newcommand{\rmap}{\mathsf{realise\text{-}map}}
\newcommand{\lp}{\mathsf{loop}}
\newcommand{\Unloop}{\mathsf{Unloop}}
\newcommand{\Const}{\mathsf{Tr}}
\newcommand{\Z}{\mathbb{Z}}
\newcommand{\cls}[1]{[#1]}
\newcommand{\cmp}{\mathbin{;}}
\title{Arbitrary pointed maps between Postnikov triples}
\date{28 October 2026}
\begin{document}
\setlength{\emergencystretch}{3em}
\maketitle


\noindent\textbf{Purpose.} The informal proof that is formalised, step by
step (each statement names its declarations; Section~9 gives the
correspondence), of the classification of \emph{arbitrary}
pointed maps between pointed simply connected $3$-types by maps of Postnikov
triples (Theorem~\ref{thm:main}), the induced equivalence of homotopy
categories (Section~\ref{sec:cat}), and the reformulation of the two sphere
arguments through it (Sections~\ref{sec:pi3}, \ref{sec:pi4}). The existing
development (\code{sag/pi3-s2-postnikov}) classifies the \emph{objects}
(\code{PostnikovClassification.classification}) and proves specialised mapping
theorems (\code{MapFib}, \code{MarkDirect}); it does not classify arbitrary
maps, and those endpoints are not used as the reason for the results below.

\noindent\textbf{Attribution.} The strategy of identifying the sphere
universal objects on the triple side, and the embedding of pointed
$2$-connected $4$-types into pointed simply connected $3$-types by looping, are
David Wärn's, as communicated by the user; the delooping machinery
(\code{DeloopMaps}, Corollary~11) is his. The object classification is
Hart--Milner's (Owen Milner's equivalence). We do not know Owen Milner's
(unpublished) treatment of maps and do not claim to reproduce it.

\noindent\textbf{Conventions.} As in the development, $\cC\,n$ means
$(n-2)$-connected and $\cT\,n$ means $(n-2)$-truncated. Composition is written
in diagrammatic order: $f\cmp g$ is ``$f$, then $g$'' (\code{ptcomp}). For
pointed $f=(f_1,f_0)$, $g$: $(f\cmp g)_1=g_1\circ f_1$ and $(f\cmp g)_0=\ap_{g_1}(f_0)\cdot
g_0$. Path composition $p\cdot q$ is \code{trans}, by $J$ on $q$, so $p\cdot
\refl\equiv p$. $K(A,n)$ is \code{EMMaps.KK[A].K n}; $K(h,n)$ is the locked
\code{H4Bilinear.Lev[A,B].K h n}. All types live in $U_l$ unless said
otherwise; groups are elements of \code{Ab} $:U_{l+1}$.

\section{The statement}

\subsection{Triples and their realisation}
A triple is $t=(A,B,k)$ with $k:K(A,2)\to_* K(B,4)$
(\code{PostnikovClassification.Triples}). Fix triples $t=(A,B,k)$,
$t'=(A',B',k')$ and write
\begin{gather*}V=K(A,2),\ Y=K(B,4),\ V'=K(A',2),\ Y'=K(B',4),\\
P_x=(k_1x=y_0),\qquad \Fib k=\textstyle\sum_{x:V}P_x,\end{gather*}
pointed by $w_0=(v_0,k_0)$, with projection $p=\fpr:\Fib k\to_* V$
(\code{PostnikovRealization.Triple.F}, \code{f0}). The facts used: $V$ is
$\cC\,3$ (Lemma~\ref{lem:unique} uses only that it is $0$-connected) and $\cT\,4$ (in fact a $2$-type); $Y$ is
$\cC\,5$; $Y'$ is $\cT\,6$; every $P_x$ is $\cC\,4$ (\code{MapFib.FibP.cp},
\code{Connected.conn-id}) and $(y=y')$ is $\cT\,5$ in $Y'$.

\subsection{Morphisms}
\begin{definition}[\code{TripleMapsEquiv.Maps2.Mor}]\label{def:mor}
\[\Mor(t,t')=\sum_{a:\Hom(A,A')}\ \sum_{b:\Hom(B,B')}
 \big(K(a,2)\cmp k'=_{\Map(V,Y')}k\cmp K(b,4)\big).\]
Accessors \code{mor-a}, \code{mor-b}, \code{mor-H}.
The third field is a \emph{path of pointed maps}: a homotopy together with
its coherence at the base point.
\end{definition}

\begin{definition}[\code{TripleMapsEquiv.Maps2.realise-map}]\label{def:rmap}
$\rmap(a,b,H)=\code{FibMap.fm}$ for the square $f=k$, $f'=k'$, $u=K(a,2)$,
$v=K(b,4)$, $sq=H$. Explicitly, with $hs_x=\code{happly}(\ap_{(-)_1}H)_x:
k'_1(u_1x)=v_1(k_1x)$ and $c=\code{pt-coh}(H):\ap_{k'_1}(u_0)\cdot k'_0=hs_{v_0}\cdot(\ap_{v_1}k_0\cdot v_0)$,
\[\rmap(a,b,H)(x,q)=(u_1x,\ hs_x\cdot(\ap_{v_1}q\cdot v_0)),\]
pointed by $\code{Sigma-path}(u_0,\ \code{fm-pt-q})$.
\end{definition}

\begin{theorem}[\code{TripleMapsEquiv.Maps2.realise-map-equiv}]\label{thm:main}
For all triples $t,t'$, $\rmap:\Mor(t,t')\to\Map(\Fib k,\Fib k')$ is an
equivalence. No sphere axiom, no hypothesis on $A,B,A',B',k,k'$.
\end{theorem}

\subsection{Sanity checks}
\emph{Levels.} $\Map(V,V')\simeq\Hom(A,A')$ and $\Map(Y,Y')\simeq\Hom(B,B')$
are sets (Lemma~\ref{lem:coef}); $\Map(V,Y')$ is $\cT\,4$ (a $2$-type) by
\code{Lemma4.cor5-map} ($V$ is $\cC\,3$, $Y'$ is $\cT\,6$), so its identity
types are $1$-types and $\Mor(t,t')$ is a $1$-type. On the other side
$\Map(\Fib k,\Fib k')$ is a $1$-type by \code{cor5-map} ($\Fib k$ is $\cC\,3$,
$\Fib k'$ is $\cT\,5$). Neither is a set in general.

\emph{The homotopy is data.} For $t=(A,0,0)$, $t'=(0,C,0)$:
$\Hom(A,0)$ and $\Hom(0,C)$ are contractible and
$\Mor(t,t')\simeq(0=_{\Map(K(A,2),K(C,4))}0)\simeq\Map(K(A,2),K(C,3))$, while
$\Fib(0)\simeq K(A,2)$ and $\Fib(0:K(0,2)\to K(C,4))\simeq\Om K(C,4)\simeq
K(C,3)$, so both sides are $\Map(K(A,2),K(C,3))$, whose components are
$H^3(K(A,2);C)$, not zero in general. The coefficients alone would give a
point. (Not formalised as a separate statement.)

\emph{Why the formula is right.} Classically it is the statement that maps of
two-stage Postnikov systems in this range are maps of bases, of fibres and a
homotopy, which follows from the Serre (or dual Blakers--Massey) estimate on
the cofibre of $\Fib k\to K(A,2)$; we do not use cofibres (no higher
inductive types), but the argument below is a direct, fibrewise replacement
for that estimate.

\section{Proof of Theorem~\ref{thm:main}}

The proof is a chain of equivalences whose composite is, by definition,
$\rmap$. Write, for a pointed $c=(c_1,c_0):V\to_* Y'$ and $Q^c_x=(c_1x=y'_0)$:
\[T(c)=\sum_{F:\prod_x P_x\to Q^c_x}\big(F_{v_0}k_0=c_0\big),\qquad
S(c)=\sum_{\beta:\Map(Y,Y')}\big(c=_{\Map(V,Y')}k\cmp\beta\big).\]
$T(c)$ is the type of pointed null-homotopies of $p\cmp c$ in fibrewise form;
$S(c)$ is the type of factorisations of $c$ through $k$.

\subsection{Maps out of $\Fib k$ into $\Fib k'$}

\begin{lemma}[\code{TripleMapsEquiv.Dec.dec-equiv}]\label{lem:dec}
The map
\begin{gather*}\mathrm{Dec}:\sum_{\alpha:\Map(V,V')}T(\alpha\cmp k')\to\Map(\Fib k,\Fib k'),\\
(\alpha,F,e)\mapsto\big((x,q)\mapsto(\alpha_1x,F_xq),\ \code{Sigma-path}(\alpha_0,\pi(e))\big)\end{gather*}
is an equivalence, where for $e:F_{v_0}k_0=\ap_{k'_1}\alpha_0\cdot k'_0$ (this
is $(\alpha\cmp k')_0$ by definition),
$\pi(e):\alpha_0{}_*(F_{v_0}k_0)=k'_0$ is
\code{transport-fib}, then $\ap_{(\ap_{k'_1}\alpha_0)^{-1}\cdot -}(e)$, then
\code{trans-cancel} (exactly the steps of \code{FibMap.fm-pt-q}).
\end{lemma}
\begin{proof}
(i) Precomposition with $p$, $(V\to V')\to(\Fib k\to V')$, is an equivalence:
\code{ConnectedMaps.Elim} at $m=\cC\,4$, along $p$ whose fibres are
$\cC\,4$ (\code{FibP.cp}), into the constant family $V'$, which is $\cT\,4$.

(ii) The map $(a_1,F)\mapsto((x,q)\mapsto(a_1x,F_xq))$,
\[\textstyle\sum_{a_1:V\to V'}\prod_{x}P_x\to(k'_1(a_1x)=y'_0)\ \to\ (\Fib k\to\Fib k'),\]
is an equivalence: it is the composite of currying
$\prod_{x}\prod_{q:P_x}\simeq\prod_{w:\Fib k}$ (inverse by
definition), \code{sigma-base} along (i), and the distributivity
$\sum_{g:\Fib k\to V'}\prod_w R(g\,w)\simeq(\Fib k\to\sum_{y}R\,y)$
(inverse by definition), for $R\,y=(k'_1y=y'_0)$.

(iii) $\Map(\Fib k,\Fib k')=\sum_{G}(G\,w_0=w'_0)$; pulling back along (ii)
(\code{sigma-base}), the fibre over $(a_1,F)$ is
$(a_1v_0,F_{v_0}k_0)=w'_0$ \emph{by definition} of the map (ii).

(iv) For fixed $(a_1,F)$, the map $(\alpha_0,e)\mapsto\code{Sigma-path}(\alpha_0,\pi(e))$,
$\sum_{\alpha_0:a_1v_0=v'_0}(F_{v_0}k_0=\ap_{k'_1}\alpha_0\cdot k'_0)\to
\big((a_1v_0,F_{v_0}k_0)=w'_0\big)$, is an equivalence: $\pi$ is a composite of
fixed paths (a left multiplication, \code{Fibres.trans-equiv}, and a right
multiplication, \code{rmul-equiv}) with $\ap$ of the equivalence
$(\ap_{k'_1}\alpha_0)^{-1}\cdot-$ (\code{ClassBase.ap-equiv}), so a fibrewise equivalence (\code{Fibres.sigma-equiv}),
followed by \code{Sigma-path}, the inverse of \code{ConnectedMaps.sigma-paths}
(\code{Whitehead.inv-equiv}).

Reordering the $\Sigma$-components ($\sum_{\alpha}\sum_F\sum_e\simeq
\sum_{a_1}\sum_F\sum_{\alpha_0}\sum_e$, inverse by definition) and composing
(ii)--(iv) gives $\mathrm{Dec}$.
\end{proof}

\subsection{The key lemma: null-homotopies are factorisations}

For $\beta:\Map(Y,Y')$ write $G^\beta_y(q)=\ap_{\beta_1}q\cdot\beta_0:\beta_1y=y'_0$
for $q:y=y_0$ (so $(k\cmp\beta)_0=G^\beta_{k_1v_0}(k_0)$), and
$Z^\beta_x=(c_1x=\beta_1(k_1x))$.

\begin{definition}[$\Phi$, \code{TripleMaps.Key.Phi}]\label{def:phi}
$\Phi_c:S(c)\to T(c)$, $\Phi_c(\beta,s)=\big(x\,q\mapsto hs_x\cdot G^\beta_{k_1x}(q),\ (\code{pt-coh}\,s)^{-1}\big)$,
with $hs=\code{happly}(\ap_{(-)_1}s)$ and
$\code{pt-coh}\,s:c_0=hs_{v_0}\cdot(k\cmp\beta)_0$.
\end{definition}

\begin{lemma}[\code{TripleMaps.Key.key}]\label{lem:key}
For every pointed $c:V\to_*Y'$, $\Phi_c$ is an equivalence.
\end{lemma}

The proof occupies the rest of this subsection. For $x:V$, $F_x:P_x\to Q^c_x$
and $\beta$, put
\[\Const^\beta_x(F_x)=\sum_{z:Z^\beta_x}\big((q\mapsto z\cdot G^\beta_{k_1x}q)=_{P_x\to Q^c_x}F_x\big)
\qquad\text{(``$F_x$ is a translate of $G^\beta$'')}.\]

\begin{lemma}[paths of pointed maps; \code{TripleMaps.PtPath.equiv}]\label{lem:ptpath}
For pointed $c,g:X\to_*Y$, $s\mapsto(\code{happly}(\ap_{(-)_1}s),\code{pt-coh}\,s)$,
\[(c=g)\to\textstyle\sum_{h:\prod_x c_1x=g_1x}(c_0=h_{x_0}\cdot g_0),\]
is an equivalence.
\end{lemma}
\begin{proof}
Fundamental theorem of identity types (\code{FundThm.FT}): the family on the
right is pointed at $g=c$ by $(\refl,\ (\code{trans-refl-l})^{-1})$, which is
the value of the map at $\refl$ (both by $J$), and its total space
$\sum_{g}\sum_h(c_0=h_{x_0}\cdot g_0)$ is contractible: $\sum_{g_1}\sum_h$ is
the singleton of $c_1$ (function extensionality), and then
$\sum_{g_0}(c_0=\refl\cdot g_0)\simeq\sum_{g_0}(c_0=g_0)$ is a singleton.
The map so obtained by $J$ agrees with
$(\code{happly}\circ\ap,\code{pt-coh})$ by $J$ on $s$.
\end{proof}

\begin{lemma}[translates; \code{TripleMaps.Transl.equiv}, \code{Transl.prop}]\label{lem:transl}
For each $x$, $z\mapsto(q\mapsto z\cdot G^\beta_{k_1x}q,\ (z,\refl))$ is an
equivalence $Z^\beta_x\simeq\sum_{F_x}\Const^\beta_x(F_x)$, and
$\Const^\beta_x(F_x)$ is a proposition.
\end{lemma}
\begin{proof}
The inverse is $(F_x,z,E)\mapsto z$; one round trip is $\refl$, the other is
$J$ on $E$. For the proposition: $\Const^\beta_x(F_x)$ is the fibre over
$F_x$ of $\tau:z\mapsto(q\mapsto z\cdot G^\beta q)$, and $\tau$ is the constant
embedding $Z^\beta_x\to(P_x\to Z^\beta_x)$ (\code{FibreH3.const-emb}: $P_x$ is
$\cC\,4$, $Z^\beta_x$ is $\cT\,5$) followed by postcomposition with the
fibrewise equivalence $d\mapsto(q\mapsto d\,q\cdot G^\beta q)$
(\code{ConnectedMaps.pi-equiv} of right multiplication by a path,
\code{TripleMaps.rmul-equiv}, by $J$ on the path); $\ap$ of the composite is
the composite of the $\ap$'s (\code{LoopCalc.ap-ap}) and $\ap$ of an equivalence
is an equivalence (\code{ClassBase.ap-equiv}); an embedding has propositional
fibres (\code{Groupoid.emb-fib-prop}).
\end{proof}

\begin{lemma}[\code{TripleMaps.Key.fac-equiv}]\label{lem:fac}
For fixed $\beta$, the map
\[(c=k\cmp\beta)\to\sum_{(F,e):T(c)}\prod_x\Const^\beta_x(F_x),\qquad
s\mapsto(\Phi_c(\beta,s),\ x\mapsto(hs_x,\refl))\]
is an equivalence.
\end{lemma}
\begin{proof}
Lemma~\ref{lem:ptpath} at $g=k\cmp\beta$ gives
$(c=k\cmp\beta)\simeq\sum_{h:\prod_xZ^\beta_x}(c_0=h_{v_0}\cdot G^\beta_{k_1v_0}k_0)$.
Lemma~\ref{lem:transl} under $\prod_x$ (\code{pi-equiv}) and distributivity
give $\prod_xZ^\beta_x\simeq\sum_F\prod_x\Const^\beta_x(F_x)$, with
$h\mapsto(x\,q\mapsto h_x\cdot G^\beta q,\ x\mapsto(h_x,\refl))$; in
particular $F_{v_0}k_0\equiv h_{v_0}\cdot G^\beta_{k_1v_0}k_0$ by
definition, so the pointing condition is carried by \code{sigma-base} to
$(c_0=F_{v_0}k_0)$, and \code{sym} turns it into $F_{v_0}k_0=c_0$. Reordering
the components gives the stated map.
\end{proof}

\begin{definition}[restriction to the fibre; \code{TripleMaps.Key.res}]
For $(F,e):T(c)$, $\mathrm{res}(F,e):\Om Y\to_*\Om Y'$ is
$q\mapsto c_0^{-1}\cdot F_{v_0}(k_0\cdot q)$, pointed by
$c_0^{-1}\cdot F_{v_0}k_0=c_0^{-1}\cdot c_0=\refl$ (via $e$; note
$k_0\cdot\refl\equiv k_0$). (This is \code{FibreH3.FibAtC.res}, conjugated by $c_0$.)
\end{definition}

\begin{lemma}[\code{TripleMaps.Key.to-loop}, \code{from-loop}]\label{lem:atpt}
For $(F,e):T(c)$ and $\beta:\Map(Y,Y')$, the propositions
$\Const^\beta_{v_0}(F_{v_0})$ and $(\mathrm{res}(F,e)=_{\Map(\Om Y,\Om Y')}\lp\,\beta)$
are logically equivalent.
\end{lemma}
\begin{proof}
The second is a proposition because $\Map(\Om Y,\Om Y')$ is a set
(\code{cor5-map}: $\Om Y$ is $\cC\,4$, $\Om Y'$ is $\cT\,5$). We use the
propositional equation
$\lp\,\beta(q)=\beta_0^{-1}\cdot(\ap_{\beta_1}q\cdot\beta_0)$
(\code{Maps.loop-ap}; it is not the definition of $\lp$).

($\Rightarrow$) Let $(z,E)$, so $F_{v_0}w=z\cdot(\ap_{\beta_1}w\cdot\beta_0)$
for all $w:P_{v_0}$ (\code{happly} $E$). Since $\Om(\Om Y')$ is $\cC\,2$,
pointed maps $\Om Y\to_*\Om Y'$ with pointwise equal underlying maps are
merely equal (\code{FibreH3.pi-mere}), and the goal is a proposition, so it
suffices that $c_0^{-1}\cdot F_{v_0}(k_0\cdot q)=\beta_0^{-1}\cdot(\ap_{\beta_1}q\cdot\beta_0)$
for all $q$. From $e$ and $E$ at $w=k_0$:
$c_0=z\cdot(\ap_{\beta_1}k_0\cdot\beta_0)$. The claim is then the groupoid
identity
\[(z\cdot(\ap_{\beta_1}k_0\cdot\beta_0))^{-1}\cdot(z\cdot(\ap_{\beta_1}(k_0\cdot q)\cdot\beta_0))
=\beta_0^{-1}\cdot(\ap_{\beta_1}q\cdot\beta_0),\]
valid for all $k_0:y=y_0$, $z$ and $q$ (\code{TripleMaps.grp1}: $(z\cdot A)^{-1}\cdot(z\cdot B)=A^{-1}\cdot B$, by
\code{idJR} on $z$; \code{Groupoid.ap-trans}; \code{TripleMaps.grp3}:
$(A\cdot b_0)^{-1}\cdot((A\cdot M)\cdot b_0)=b_0^{-1}\cdot(M\cdot b_0)$, by $J$ on
$b_0$).

($\Leftarrow$) Let $r:\mathrm{res}(F,e)=\lp\,\beta$. Take
$z=F_{v_0}k_0\cdot(\ap_{\beta_1}k_0\cdot\beta_0)^{-1}$. For $w:P_{v_0}$ put
$q=k_0^{-1}\cdot w$, so $w=k_0\cdot q$ (\code{trans-cancel'}); then
$F_{v_0}w=c_0\cdot\mathrm{res}(F,e)(q)=c_0\cdot\beta_0^{-1}\cdot\ap_{\beta_1}q\cdot\beta_0$
(\code{happly} $r$), and with $e:F_{v_0}k_0=c_0$ the required
$z\cdot(\ap_{\beta_1}w\cdot\beta_0)=F_{v_0}w$ follows from \code{ap-trans}, \code{tassoc}, \code{grp3} and
\code{trans-cancel'} (\code{TripleMaps.Key.goal-at}); it is proved at $w=k_0\cdot q$
and transported along $k_0\cdot(k_0^{-1}\cdot w)=w$. Function extensionality
gives $E$.
\end{proof}

\begin{lemma}[\code{TripleMaps.Key.unique-fac}]\label{lem:unique}
For every $(F,e):T(c)$, $\sum_{\beta}\prod_x\Const^\beta_x(F_x)$ is contractible.
\end{lemma}
\begin{proof}
For each $\beta$, $\prod_x\Const^\beta_x(F_x)$ and
$\Const^\beta_{v_0}(F_{v_0})$ are logically equivalent propositions: one
direction is evaluation, the other is \code{conn-elim} over the $0$-connected
$V$ into the propositions $\Const^\beta_x(F_x)$ (Lemma~\ref{lem:transl}; this
is \code{FibreH3.spread}, there stated at $V=K(\Z,2)$). With
Lemma~\ref{lem:atpt}, $\sum_\beta\prod_x\Const^\beta_x(F_x)\simeq
\sum_\beta(\mathrm{res}(F,e)=\lp\,\beta)$ (\code{Logic.sigma-equiv} at the
identity), and the latter is $\Unloop$ of $\mathrm{res}(F,e)$, contractible by
\code{DeloopMaps.Maps.Over.cor11} at $c=2$ ($Y$ is $\cC\,4$, $Y'$ is
$\cT\,6$), as in \code{FibreH3.FibAtC.deloop-c}.
\end{proof}

\begin{proof}[Proof of Lemma~\ref{lem:key}]
$\Phi_c$ is the composite
\[S(c)\xrightarrow{\ \sum_\beta\text{Lemma~\ref{lem:fac}}\ }
\sum_\beta\sum_{(F,e)}\prod_x\Const^\beta_x(F_x)
\xrightarrow{\ \text{swap}\ }\sum_{(F,e)}\sum_\beta\prod_x\Const^\beta_x(F_x)
\xrightarrow{\ \mathrm{pr}_1\ }T(c),\]
the swap has inverse by definition, and $\mathrm{pr}_1$ is an equivalence by
Lemma~\ref{lem:unique} (\code{SphereTruncation.fst-equiv}). The composite is
$\Phi_c$ by definition.
\end{proof}

\begin{remark}
This is where the translation data live: a fibrewise map
$F_x:P_x\to Q^c_x$ between two ``$K(B,3)$-torsors'' is a linear part (the
delooping $\beta$ of its restriction to the fibre over $v_0$, whose
coefficient is $b$) and a translation (the constant $z=h_x$, the homotopy).
The base point fixes the translation at $v_0$ through $e$; no point of any
other fibre and no representative of a class is chosen. Lifts are not
assumed unique: the ambiguity is exactly the homotopy component $h$, which is
kept as data.
\end{remark}

\subsection{Coefficients}
\begin{lemma}[\code{TripleMapsEquiv.Coef.k2-equiv}, \code{k4-equiv}]\label{lem:coef}
$a\mapsto K(a,2):\Hom(A,A')\to\Map(V,V')$ and
$b\mapsto K(b,4):\Hom(B,B')\to\Map(Y,Y')$ are equivalences.
\end{lemma}
\begin{proof}
Level $2$: inverse \code{H4Bilinear.H2.inv2}, with \code{k2-surj} and, for
$\mathsf{inv2}(K(a,2))=a$, \code{k2-inj} applied to \code{k2-surj}. Level $4$:
inverse \code{FibreH3.Climb.inv4} with \code{surj4}; for
$\mathsf{inv4}(K(b,4))=b$: $\lp\,K(b,m+1)=\mathsf{conj}_m K(b,m)$
(\code{Lev.loopK}) and $\mathsf{unconj}_m(\mathsf{conj}_mF)=F$
(\code{TripleMapsEquiv.Coef.unconj-conj}, the computation of \code{H2.unconj-conj} at
level $m$: $e_A^{-1}\cmp e_A\cmp F\cmp e_B^{-1}\cmp e_B=F$), so
$\mathsf{inv4}(K(b,4))=\mathsf{inv3}(K(b,3))=\mathsf{inv2}(K(b,2))=b$.
\end{proof}

\subsection{Assembly}
$\rmap$ is the composite
\begin{align*}
\Mor(t,t')&\xrightarrow{\ (a,b,H)\mapsto(a,K(b,4),H)\ }\textstyle\sum_{a}S(K(a,2)\cmp k')
\xrightarrow{\ \sum_a\Phi\ }\textstyle\sum_aT(K(a,2)\cmp k')\\
&\xrightarrow{\ (a,\theta)\mapsto(K(a,2),\theta)\ }\textstyle\sum_{\alpha}T(\alpha\cmp k')
\xrightarrow{\ \mathrm{Dec}\ }\Map(\Fib k,\Fib k'),
\end{align*}
the first and third by \code{sigma-base} along Lemma~\ref{lem:coef}, the
second by Lemma~\ref{lem:key} and \code{Fibres.sigma-equiv}. That the composite
\emph{is} \code{FibMap.fm} holds by definition: the function is
$(x,q)\mapsto(u_1x,hs_x\cdot(\ap_{v_1}q\cdot v_0))$ and the pointing is
$\code{Sigma-path}(u_0,\pi((\code{pt-coh}H)^{-1}))$, which is
\code{fm-pt} ($c\equiv\code{pt-coh}\,H$, $\pi$ is the calc of \code{fm-pt-q}).
This is recorded as a \code{refl} lemma \code{realise-map-def} (rule~3), and
Theorem~\ref{thm:main} follows.

\subsection{What the canonical map does}
(a) $\rmap(a,b,H)\cmp p'=p\cmp K(a,2)$: \code{FibMap.fm-over}. Hence on $\pi_2$
(through $p$, $p'$) it induces $a$.
(b) On the fibre over the base point it is $\Phi$'s $F_{v_0}$, and
$\mathrm{res}(\Phi(K(b,4),H))=\lp\,K(b,4)$ (Lemma~\ref{lem:atpt}
($\Rightarrow$) with $z=hs_{v_0}$): the restriction to the fibres is the looping of
$K(b,4)$, which is $\mathsf{conj}_3K(b,3)$ (\code{loopK}); this is how the
coefficient $b$ is read off on $\pi_3$.
(c) Conversely, the inverse sends $G$ to $a=\mathsf{inv2}$ of the
extension of $G\cmp p'$ along $p$ (Lemma~\ref{lem:dec}(i)) and $b$ to the
coefficient of the delooping of the restriction (Lemma~\ref{lem:unique}).

\section{Identities, composition, homotopy categories}\label{sec:cat}

\begin{definition}[\code{TripleHomotopy.Maps3.mcomp}, \code{Maps1.mid}]
For $m=(a,b,H):t\to t'$, $m'=(a',b',H'):t'\to t''$:
$m\cmp m'=(a\cmp a',\ b\cmp b',\ H'')$ with $H''$ the composite
\[K(a\cmp a',2)\cmp k''=K(a,2)\cmp K(a',2)\cmp k''=K(a,2)\cmp k'\cmp K(b',4)
=k\cmp K(b,4)\cmp K(b',4)=k\cmp K(b\cmp b',4),\]
i.e.\ $\ap_{-\cmp k''}$ of \code{kc-lev} at $2$ (\code{KFun.Comp.k-comp} through the
locks \code{Lev.K-def}), then the pasting \code{Paste.paste} (\code{ptcomp-assoc},
$\ap_{K(a,2)\cmp-}H'$, \code{ptcomp-assoc}$^{-1}$, $\ap_{-\cmp K(b',4)}H$,
\code{ptcomp-assoc}), then $\ap_{k\cmp-}$ of \code{kc-lev}$^{-1}$ at $4$
(\code{Bridge.bridge}). $\mathrm{id}_t=(1_A,1_B,H_1)$, $H_1$ the bridge along
$K(1,2)=\mathrm{id}$ and $K(1,4)=\mathrm{id}$ (\code{KId.k2-id}, \code{k4-id}) of
\code{ptcomp-idl}$\,\cdot\,$\code{ptcomp-idr}$^{-1}$ (\code{FmId.idsq}).
$K(1,m)=\mathrm{id}$ for $m=2,4$ is proved from the surjectivity of $K(-,m)$: if
$K(h,m)=\mathrm{id}$ then $K(1,m)=K(1,m)\cmp K(h,m)=K(1\cmp h,m)=K(h,m)=\mathrm{id}$.
\end{definition}

\begin{proposition}[\code{rmap-comp-pt}, \code{rmap-comp}, \code{rmap-id-pt}, \code{rmap-id}]\label{prop:funct}
$\rmap(\mathrm{id}_t)$ and the pointed identity, and $\rmap(m\cmp m')$ and
$\rmap(m)\cmp\rmap(m')$, have pointwise equal underlying functions; hence
they are merely equal as pointed maps (\code{pi-mere}, $\Om\Fib k''$ being
$\cC\,2$), i.e.\ equal in $[\Fib k,\Fib k'']_*$.
\end{proposition}
\begin{proof}
Write $hp(r)_x=\code{happly}(\ap_{(-)_1}r)_x$ for the underlying homotopy of a path
$r$ of pointed maps (\code{hp}). By $J$: $hp(r\cdot s)=hp(r)\cdot hp(s)$,
$hp(r^{-1})=hp(r)^{-1}$, $hp(\ap_{u\cmp-}s)_x=hp(s)_{u_1x}$,
$hp(\ap_{-\cmp v}s)_x=\ap_{v_1}(hp(s)_x)$; $hp(\code{ptcomp-assoc})=\refl$ ($J$ on
the three pointings, as in its definition); a path changing only the pointing,
such as \code{ptcomp-idl}, \code{ptcomp-idr}, has trivial homotopy
(\code{hp-pair}).
(1) \code{Bridge.fm-bridge}: for paths $U:u_2=u$, $W:v_2=v$ of pointed maps and a
square $s$ for $(u,v)$, \code{fm} of the bridged square for $(u_2,v_2)$ is pointwise
\code{fm} of $s$ ($J$ on $U$, $W$, whose left endpoints are free; at $\refl$ the
bridge is $\refl\cdot s$, \code{trans-refl-l}).
(2) \code{Paste.fm-paste}: the first components agree by definition; the
homotopy of the pasting is $hs'_{u_1x}\cdot\ap_{v'_1}(hs_x)$ (\code{hs-paste},
the rules above), and the second components agree by \code{ap-trans},
\code{ap-ap} and the groupoid identity \code{grp4} (\code{eq2}).
(3) \code{FmId.fm-id}: the homotopy of \code{idsq} is $\refl$, and
$\refl\cdot\ap_{\mathrm{id}}q=q$.
Then \code{rmap-comp-pt} is (1) followed by (2), \code{rmap-id-pt} is (1) followed
by (3).
\end{proof}

\noindent A coherent pointed path (the base-point coherence) is \emph{not}
constructed; the homotopy category only needs classes (as confirmed by the
user, this suffices for David Wärn's argument).

\begin{theorem}[\code{TripleHomotopy.Cat}]\label{thm:hocat}
Put $\hom(t,t')=\|\Mor(t,t')\|_0$ (\code{homT}) and $R=\|\rmap\|_0:\hom(t,t')\to
[\Fib k,\Fib k']_*$ (\code{R}), a bijection by Theorem~\ref{thm:main} and
\code{classes-equiv} (\code{R-equiv}). Composition and identities on $\hom$ are
transported from \code{compC}, \code{idC} along $R$ (\code{compT}, \code{idT};
$R$ preserves them: \code{R-comp}, \code{R-id}), and they \emph{are} the intrinsic
ones on representatives: $\cls m\cmp\cls{m'}=\cls{m\cmp m'}$ and
$\mathrm{id}=\cls{\mathrm{id}_t}$ (\code{compT-rep}, \code{idT-rep}, by
Proposition~\ref{prop:funct} and injectivity of $R$). The category laws hold
(\code{compT-assoc}, \code{compT-idl}, \code{compT-idr}, from \code{comp-assoc},
\code{comp-idl}, \code{comp-idr}). So $(\real,R)$ is a functor from the triple
homotopy category to the homotopy category of pointed simply connected $3$-types
which is an equivalence on object types
(\code{PostnikovClassification.classification}; object types are not sets, and
no univalence of either category is claimed) and bijective on hom-sets: an
equivalence of homotopy categories, with inverse $\real^{-1}$ on objects and
$R^{-1}$ composed with transport along $\real(\real^{-1}X)=X$ on hom-sets (the
latter is stated here, not formalised as a separate declaration).
\end{theorem}

\section{The components of the compatible homotopies}\label{sec:comp}

Both sphere arguments compute $\hom(t_0,t')$ directly from
$\Mor(t_0,t')=\sum_{a,b}(K(a,2)\cmp k'=\kappa_0\cmp K(b,4))$, $\kappa_0$ the source
invariant: the coefficient $b$ is determined by $a$ through (H4) (with
Statement~T for $\bar k$), and the components of the identity type are
controlled by (H3). The latter needs loops at an arbitrary point of
$M=\Map(K(\Z,2),Y')$, $Y'=K(B',4)$; we reduce them to loops at the base
point by translation.

\begin{lemma}[translation; \code{TripleComponents.PtTransl}]\label{lem:translation}
Let $X$, $Z$ be pointed and $M'=\Map(X,\Om Z)$. For $G:M'$, the map
$\mathrm{tr}_G(F)=(x\mapsto F_1x\cdot(G_1x)^{-1},\ \text{pointing from }F_0,G_0)$
is an equivalence $M'\to M'$ (inverse: right multiplication by $G$), and
$\mathrm{tr}_G(G)$ is merely equal to the constant map $0$ when $\Om(\Om Z)$ is
$\cC\,2$ (\code{pi-mere}; pointwise \code{trans-sym-r}).
\end{lemma}

\begin{lemma}[\code{TripleComponents.Unique.unique}, \code{unique-pt}]\label{lem:components}
Let $X$ be pointed, $Y'$ pointed with a pointed equivalence
$e:\Om Z\simeq_*Y'$ ($Z=K(B',5)$, $e=e_4$), $\Om^2Z$ $\cC\,2$, and assume every
pointed map $X\to_*\Om^2 Z$ is merely constant. Then for $f,g:\Map(X,Y')$ and
$H,H':f=g$, $H$ and $H'$ are merely equal.
\end{lemma}
\begin{proof}
Postcomposition with $e^{-1}$ is an equivalence $\Map(X,Y')\simeq M'$, so $\ap$
of it is injective; then $\ap\,\mathrm{tr}_{e^{-1}g}$ (Lemma~\ref{lem:translation})
reduces to two paths $r,r':u=v$ in $M'$ with $v$ merely $0$; the goal being a
proposition, $v=0$ by $J$, and $r^{-1}\cdot r'$ is a loop at $0$ in $M'$,
i.e.\ (\code{CupSquare.PtMapLoops}) a pointed map $X\to_*\Om^2Z$, merely
constant by hypothesis; so $r^{-1}\cdot r'$ is merely $\refl$.
\end{proof}

\noindent For $X=K(\Z,2)$: $\Om^2K(B',5)\simeq_*K(B',3)$ ($\lp\,e_4$, $e_3$) and
(H3) (\code{H3Proof.h3}: $[K(\Z,2),K(B',3)]_*$ is a point) give the hypothesis.

\section{$\pi_3(S^2)$ through triple morphisms}\label{sec:pi3}

Let $t_0=(\Z,\Z,\kappa)$, $P=\Fib\kappa$ (\code{CupSquare.kappa}), and
$t'=(A',B',k')$ arbitrary.

\begin{proposition}[\code{Pi3TripleUniversal.ThetaGen.theta-equiv}, \code{Theta}]\label{prop:theta}
$\theta:\hom(t_0,t')\to\Hom(\Z,A')$, $\cls{m}\mapsto\code{mor-a}\,m$, is a
bijection.
\end{proposition}
\begin{proof}
$\Hom(\Z,A')$ is a set, so $\theta$ is well defined.
\emph{Surjective:} for $a$, (H4) (\code{H4Proof.h4}: $b\mapsto\cls{\kappa\cmp
K(b,4)}$ is a bijection $\Hom(\Z,B')\to[K(\Z,2),Y']_*$) gives $b$ with
$\cls{\kappa\cmp K(b,4)}=\cls{K(a,2)\cmp k'}$, hence merely a path $H$ and an
element $(a,b,H^{-1})$ over $a$. \emph{Injective:} if $\code{mor-a}\,m=\code{mor-a}\,m'=a$
then $\cls{\kappa\cmp K(b,4)}=\cls{K(a,2)\cmp k'}=\cls{\kappa\cmp K(b',4)}$, so
$b=b'$ by (H4); transporting, $H,H':K(a,2)\cmp k'=\kappa\cmp K(b,4)$, merely
equal by Lemma~\ref{lem:components} with (H3); so $\cls m=\cls{m'}$.
\end{proof}

\begin{proposition}[\code{Pi3TripleUniversal.Final.At.marked}]\label{prop:mark3}
Let $s=\code{MapFib.Final.s}\in[S^2,P]_*$ ($s\cmp\cls p=\cls\sigma$,
\code{s-eq}). For every triple $t'$, $s^*:[P,\Fib k']_*\to[S^2,\Fib k']_*$ is a
bijection.
\end{proposition}
\begin{proof}
The composite
\[\hom(t_0,t')\xrightarrow{\|\rmap\|_0}[P,\Fib k']_*\xrightarrow{s^*}[S^2,\Fib k']_*
\xrightarrow{p'_*}[S^2,V']_*\]
sends $\cls m$ to $\cls{\sigma\cmp K(a,2)}$: $s\cmp\cls{\rmap\,m}\cmp\cls{p'}=s\cmp\cls{p}\cmp\cls{K(a,2)}$
(\code{fm-over}, \code{comp-assoc}) $=\cls\sigma\cmp\cls{K(a,2)}$
(\code{s-eq}). The map $a\mapsto\cls{\sigma\cmp K(a,2)}$, $\Hom(\Z,A')\to[S^2,V']_*$,
is a bijection: $\sigma^*$ on classes into the $2$-type $V'$
(\code{MapFib.SourceS.Sig.sigma-cls}) after
$\|K(-,2)\|_0:\Hom(\Z,A')\simeq[K(\Z,2),V']_*$ (Lemma~\ref{lem:coef},
\code{classes-equiv}). So the composite is this bijection after $\theta$;
$\|\rmap\|_0$ (Theorem~\ref{thm:main}, \code{classes-equiv}) and $p'_*$ (the
lifting lemma at the source $S^2$, \code{MapFib.SourceS.qS}, which uses only
\code{s2-paths}) are bijections, hence so is $s^*$ (2-out-of-3).
\end{proof}

\noindent Then for $X:\SCt$, $X=\real(\real^{-1}X)$ (\code{classification}),
$s^*$ is bijective at $X$ (transport), which is \code{Pi3Sphere}'s \code{MapThm};
\code{Pi3Sphere.Conditional.pi3-S2-T} (truncation, two-map Yoneda,
\code{HomotopyGroups}) gives $\pi_3(S^2)\cong\Z$:
\code{Pi3TripleUniversal.Final.pi3-S2}. \emph{Not used:} \code{MapFib.Final.rep-fib},
\code{MapFib.Lifting} at the source $P$, \code{FibreH3.fib-h3}, \code{Presentation}.
\emph{Used and disclosed:} the marked class \code{MapFib.Final.s} (from
\code{qS} at $S^2$ and $\sigma$, not from \code{rep-fib}).

\section{$\pi_4(S^3)$ through triple morphisms and the looping embedding}\label{sec:pi4}

$\SCf$ embeds in $\SCt$ by $\Om$ (\code{DeloopEmbedding.Emb.emb},
\code{sc4-equiv}); on maps, looping is a bijection
$[X,Y]_*\simeq[\Om X,\Om Y]_*$ for $X$ $\cC\,4$, $Y$ $\cT\,6$
(\code{MarkLoop.LoopCls.equiv}), and statement M at $u:\SCf$ is equivalent to
bijectivity of $r^*:[R,\Om Y]_*\to[\Om S^3,\Om Y]_*$ with
$r=\lp\,s\cmp\cls{\omega_{BR}}$ (\code{MarkLoop.LMark.r}, Proposition L:
\code{LMark.At.from-loop}). Here $\bar t=(\Z,\Z/2,\bar k)$, $R=\Fib\bar k=\real\bar t$,
$\bar p:R\to K(\Z,2)$.

\begin{proposition}[\code{Pi4TripleUniversal.Theta4.TG.theta-equiv}]\label{prop:theta4}
For every deloopable $t'=(A',B',k')$ (\code{DeloopEmbedding.Deloopable}),
$\theta:\hom(\bar t,t')\to\Hom(\Z,A')$, $\cls m\mapsto\code{mor-a}\,m$, is a
bijection.
\end{proposition}
\begin{proof}
\emph{Surjective:} by (H4), $\cls{K(a,2)\cmp k'}=\Phi_{B'}(h)$ for some $h$; by
Statement~T (\code{TwoTorsion.TwoTors.Two.two-tors} with $G=K(a,2)$, the only
use of deloopability) $2h(1)=0$, so $h=\rho\cmp g$ (\code{Necessity.Factor}) and
$\cls{K(a,2)\cmp k'}=\Psi(g)=\cls{\bar k\cmp K(g,4)}$ (\code{k-comp},
\code{ptcomp-assoc}); so merely $(a,g,H)$.
\emph{Injective:} for $m,m'$ over $a$, $\Psi(b)=\cls{K(a,2)\cmp k'}=\Psi(b')$, so
$b=b'$ (\code{Necessity.Psi}, injective); the homotopies are merely equal by
Lemma~\ref{lem:components} with (H3).
\end{proof}

\begin{proposition}[\code{Pi4TripleUniversal.Final4.marked}]\label{prop:mark4}
For every $u:\SCf$, M holds at $u$.
\end{proposition}
\begin{proof}
Let $t'=\real^{-1}(\Om u)$, deloopable by $u$, $\varepsilon:\Om Y\simeq_*\Fib k'$.
\emph{Lifting at the source $\Om S^3$, without} \code{MarkDirect.ExR}: we use
\code{MapFib.Lifting} at $Q=\Om S^3$ with hypothesis (U) from
\code{MarkDirect.Direct.UnS} (sphere side) and hypothesis (E) proved anew:
for $F:\Om S^3\to_*V'$, by $r^*$ bijective into $K(A',2)$
(\code{MarkDirect.Direct.rA}), $\bar p^*$ bijective (\code{MapFib.FibP[}$\Z,\Z/2,\bar k$\code{].PStar})
and \code{k2-surj}, $\cls F=r\cmp\cls{\bar p\cmp K(a,2)}$ for some $a$; by
Proposition~\ref{prop:theta4} (surjectivity) there is merely $(a,b,H)$, and
then $F\cmp k'$ is merely $r\cmp\bar p\cmp\bar k\cmp K(b,4)$, merely $0$
(\code{FibP.fseq} at $\bar t$). Hence $p'_*:[\Om S^3,\Fib k']_*\to[\Om
S^3,V']_*$ is bijective. Then the composite
\[\hom(\bar t,t')\xrightarrow{\|\rmap\|_0}[R,\Fib k']_*\xrightarrow{r^*}[\Om S^3,\Fib k']_*
\xrightarrow{p'_*}[\Om S^3,V']_*\]
sends $\cls m$ to $r\cmp\cls{\bar p}\cmp\cls{K(a,2)}$ (\code{fm-over}), and
$a\mapsto r\cmp\cls{\bar p\cmp K(a,2)}$ is a bijection
$\Hom(\Z,A')\to[\Om S^3,V']_*$ (Lemma~\ref{lem:coef}, \code{PStar},
\code{rA}). So $r^*$ is bijective at $\Fib k'$ (2-out-of-3), at $\Om Y$
(\code{MarkDirect.Direct.transfer} along $\varepsilon^{-1}$), and M holds at $u$
(Proposition L).
\end{proof}

\noindent Then $\pi_4(S^3)\cong\Z/2$ as in \code{Pi4Direct.Assemble}, with
\emph{both} instances $s_P$ ($u=BR$) and $s_T$ ($u=\|S^3\|_4$) from this
proposition, \code{Pi4Sphere.Two4} and \code{Pi4Sphere.pi4-BR}:
\code{Pi4TripleUniversal.Final4.pi4-S3}.
\begin{itemize}\raggedright
\item \emph{Not used:} \code{MarkDirect.Direct.rF}, \code{marked}, \code{qS},
\code{ExS}; \code{MarkDirect.ExR}, \code{qR}; \code{DeloopEmbedding.Mark.marked};
\code{SourceS3.Final.rep-fib}; \code{Presentation4}; \code{MapBR.ur}.
\item \emph{Used and disclosed:} the marked classes \code{SourceS3.Final.s} and
\code{MarkLoop.LMark.r} (from \code{qS3} and $\sigma_3$);
\code{MarkDirect.Direct.UnS}; \code{MarkDirect.Direct.rA} (through
\code{SourceS3.Final.right}: $\sigma_3^*$ on classes into $K(A,3)$);
\code{MarkDirect.Direct.transfer}.
\end{itemize}

\section{Universes}
For fixed groups, all of $\Mor(t,t')$, $T$, $S$, $\Map(\Fib k,\Fib k')$ and
$\hom(t,t')$ live in $U_l$; the types of objects $\Tri$, $\SCt$, $\SCf$ live in
$U_{l+1}$. No universe change is needed in the new modules. In particular, the
lemma \code{ClassBase.conn-lift}, which turns a connectivity statement made with
truncations in $U_l$ into the same statement in $U_{l+1}$ and which the object
classification needs, is not used: Theorem~\ref{thm:main} is stated for fixed
$t,t':\Tri$, and in Sections~\ref{sec:pi3}--\ref{sec:pi4} the transport along
$\real(\real^{-1}X)=X$ is a path in $\SCt:U_{l+1}$ acting on statements in
$U_l$, as in \code{Pi3Classified}.

\section{Declaration correspondence}
\begin{tabular}{>{\raggedright}p{0.36\linewidth}>{\raggedright\arraybackslash}p{0.58\linewidth}}
Definitions~\ref{def:mor}, \ref{def:rmap} & \code{TripleMapsEquiv.Maps2.Mor}, \code{mor-a/b/H}, \code{ck}, \code{kc}, \code{realise-map}, \code{realise-map-def} (\code{refl})\\
$T(c)$, Definition~\ref{def:phi} & \code{TripleMaps.Tc}; \code{TripleMaps.Key.S}, \code{T}, \code{G}, \code{hs}, \code{Phi}\\
Lemma~\ref{lem:ptpath} & \code{TripleMaps.PtR}, \code{ptpath}, \code{PtPath.equiv}\\
Lemma~\ref{lem:transl} & \code{TripleMaps.rmul-equiv}, \code{Transl.equiv}, \code{Transl.prop}; \code{Key.Tl}, \code{tl-prop}\\
Lemma~\ref{lem:fac} & \code{Key.m2-equiv}, \code{fac-map}, \code{fac-equiv}\\
Lemma~\ref{lem:atpt} & \code{Key.res}, \code{grp1}, \code{grp3}, \code{to-loop-at}, \code{to-loop}, \code{zt}, \code{goal-at}, \code{from-loop}; \code{res-phi}\\
Lemma~\ref{lem:unique} & \code{Key.unique-fac}\\
Lemma~\ref{lem:key} & \code{TripleMaps.Key.key}\\
Lemma~\ref{lem:dec} & \code{TripleMapsEquiv.Dec.piq}, \code{dec}, \code{piq-equiv}, \code{cur-equiv}, \code{bmap-equiv}, \code{dec-equiv}\\
Lemma~\ref{lem:coef} & \code{TripleMapsEquiv.Coef.unconj-conj}, \code{inv2-K}, \code{inv3-K}, \code{inv4-K}, \code{k2-equiv}, \code{k4-equiv}\\
Theorem~\ref{thm:main} & \code{TripleMapsEquiv.Maps2.realise-map-equiv}\\
Section~\ref{sec:cat} & \code{TripleHomotopy}: \code{hp*}, \code{Bridge}, \code{Paste}, \code{FmId}, \code{kc-lev}, \code{KId}, \code{Maps3}, \code{Maps1}, \code{Cat}\\
Lemmas~\ref{lem:translation}, \ref{lem:components} & \code{TripleComponents.PtTransl}, \code{Unique}, \code{unique-pt}, \code{const-pt}\\
Section~\ref{sec:pi3} & \code{Pi3TripleUniversal.H3U}, \code{KCls}, \code{ThetaGen} (Proposition~\ref{prop:theta} for any source triple $(\Z,C,k_s)$ from existence and uniqueness of $b$), \code{Theta}, \code{Final.At.marked}, \code{Final.pi3-S2}\\
Section~\ref{sec:pi4} & \code{Pi4TripleUniversal.Theta4}, \code{Final4.At} (\code{chain}, \code{null1}, \code{ExS4}, \code{qS4}, \code{rr-equiv}, \code{rFt}), \code{Final4.AtU}, \code{marked}, \code{Assemble}, \code{Final4.pi4-S3}\\
\end{tabular}

\section{Validation}
Fresh checks (28 October 2026): \code{./segc sag/pi3-s2-postnikov all}: 102
modules, 0 failed; a new process of \code{cabal list-bin exe:seg} on
\code{All.seg}: 5633 declarations, 0 failed, 3.4\,s. The uniform checker
(\code{seg/headlinear}, \code{segd --opens}) passes all 32 \code{.opens}
specifications (the 26 existing ones and \code{triplemaps}, \code{triplemapsequiv},
\code{triplecomponents}, \code{triplehomotopy}, \code{pi3triple}, \code{pi4triple}), every
entry ``nothing else opened''. The six new specifications were generated from
\code{--opened} and audited against this note by kind (listed in their headers),
not entry by entry; no construction of $K(h,m)$, $e_m$, the deloopings or the cup
square is opened by the new declarations.

\emph{Dependency test.} In a scratch copy (not the working tree) the old
specialised results were deleted: \code{MapFib.Final.rep-fib}, \code{rep-pteq},
\code{rep-cup}, \code{SourceP.qP}, \code{ExP}, \code{ex-P};
\code{MarkDirect.ExR}, \code{ex-R}, \code{un-R}, \code{qR}, \code{Direct.ExS}, \code{ex-S}, \code{qS},
\code{rF}, \code{AtU}, \code{marked}; \code{DeloopEmbedding.Mark.At}, \code{Mark.marked};
\code{SourceS3.Final.rep-fib}, \code{rep-pteq}; \code{FibreH3.fib-h3}, \code{FibAt};
\code{Presentation.Final}, \code{Presentation4.Final4}, \code{Pi4Sphere.Assemble4},
\code{Pi4Sphere.Pi4S3}, and the files \code{Pi3Classified}, \code{Pi4Direct}. The copy
checks with all new modules (5300 declarations, 0 failed). The old endpoints are
unchanged in the working tree and still check.

\section{What is not claimed}
A coherent pointed path for Proposition~\ref{prop:funct} (only equality in
$[\Fib k,\Fib k'']_*$); any bicategorical coherence; univalence of the
categories; an explicit inverse functor as a declaration (only the object
equivalence and the hom-set bijections with functoriality); a separate formal
treatment of the example $t=(A,0,0)$, $t'=(0,C,0)$; a formal statement that
$\rmap$ induces $b$ on $\pi_3$ through \code{Triple.pi3-iso} (what is formalised
is \code{fm-over}, which gives $a$ on the base, and \code{Key.res-phi}: the
restriction to the fibre over the base point is $\lp\,K(b,4)$).
\section{History}
A first version was written before any code
(\code{seg/WRITING.md}, rule~1) and refereed independently against the library.
It was then revised so that it is the proof checked. The changes, made either
after review or during formalisation, are:
the $\pi_4$ lifting at $\Om S^3$ (existence now from the surjectivity of
$\theta$, instead of \code{MarkDirect.Direct.qS}, which depended on
\code{MarkDirect.ExR});
the direct computation of the hom-sets from $\Mor$, with the translation and
components lemmas (Section~\ref{sec:comp}) replacing an argument through
\code{FibreH3.fib-h3} and \code{MapBR.ur};
right multiplication (\code{rmul-equiv}) and \code{ClassBase.ap-equiv} in
Lemmas~\ref{lem:transl} and~\ref{lem:dec};
the groupoid steps of Lemma~\ref{lem:atpt} ($\Leftarrow$), now \code{goal-at};
$K(1,m)=\mathrm{id}$ from the surjectivity of $K(-,m)$ instead of an induction;
composition on $\hom$ transported along $R$, with its agreement with the
intrinsic pasting proved (\code{compT-rep}).

\end{document}
